% ============================================================
% ABSTRACT
% University of Turku, Faculty of Technology
% MSc Thesis - Microbiome-Based BMI Prediction
% ============================================================

\begin{abstract}

The human gut microbiome has emerged as a key modulator of host metabolism and body composition. This thesis investigates whether gut microbial composition, assessed through shotgun metagenomics, can predict Body Mass Index (BMI) using supervised machine learning. A reproducible Nextflow pipeline was developed and applied to 17,903 human samples from the Metalog consortium—one of the largest metagenomic cohorts assembled for this purpose.

Four machine learning algorithms were benchmarked: Elastic Net (linear baseline), Random Forest, XGBoost, and Support Vector Machine. Random Forest achieved the highest predictive performance ($R^2 = 0.325$), more than doubling the linear baseline ($R^2 = 0.14$), demonstrating that the microbiome-BMI relationship is fundamentally non-linear. Feature importance analysis identified \emph{Adlercreutzia equolifaciens}, \emph{Roseburia lenta}, and several uncharacterized species as leading predictors, offering candidates for mechanistic investigation.

A 1\% prevalence filter reduced dimensionality by 83\% (10,701 to 1,799 features) with negligible accuracy loss, dramatically improving computational efficiency. Saturation analysis revealed diminishing returns beyond 10,000 samples, providing practical guidance for resource-constrained studies. The pipeline, trained exclusively on microbial features without demographic covariates, establishes a microbiome-only baseline for BMI prediction.

These findings confirm that the gut ecosystem encodes meaningful metabolic information accessible through standard machine learning pipelines on consumer-grade hardware. The reproducible workflow and identified taxa provide a foundation for translating microbial insights into precision nutrition and personalized medicine.

\end{abstract}


